\documentclass[11pt,a4paper]{article}
\usepackage[utf8]{inputenc}
\usepackage[T1]{fontenc}
\usepackage[ngerman]{babel}
\usepackage[margin=2cm]{geometry}
\usepackage{enumitem}
\usepackage{xcolor}
\usepackage{titlesec}
\usepackage{fancyhdr}
\usepackage{siunitx}

% Page style
\pagestyle{fancy}
\fancyhf{}
\fancyhead[L]{\small PhD-Verteidigung - Stichpunkte}
\fancyhead[R]{\small Seite \thepage}
\renewcommand{\headrulewidth}{0.4pt}

% Section formatting
\titleformat{\section}{\Large\bfseries\color{blue!70!black}}{\thesection}{1em}{}
\titleformat{\subsection}{\large\bfseries\color{blue!50!black}}{\thesubsection}{1em}{}
\titleformat{\subsubsection}{\normalsize\bfseries}{\thesubsubsection}{1em}{}

% Compact lists
\setlist[itemize]{leftmargin=*,topsep=2pt,itemsep=1pt,parsep=0pt}

% Title
\title{\textbf{PhD-Verteidigung Präsentation}\\[0.5em]
\Large Bolometrie-Diagnostik und Echtzeit-Strahlungsrückkopplungskontrolle\\
am Wendelstein 7-X\\[1em]
\large Stichpunkte Leitfaden}
\author{}
\date{\today}

\begin{document}

\maketitle
\thispagestyle{empty}

\vspace{1em}
\noindent\textbf{Anleitung:} Dieses Dokument enthält alle Stichpunkte organisiert nach Präsentationsabschnitten. Nutzen Sie es als Referenzleitfaden während der Verteidigung, um umfassende Abdeckung aller Schlüsselpunkte sicherzustellen.

\newpage
\setcounter{page}{1}

\section{TITELFOLIE}

\begin{itemize}
\item Willkommen zu meiner PhD-Verteidigung über Bolometrie-Diagnostik und Echtzeit-Strahlungsrückkopplungskontrolle am Wendelstein 7-X
\item Diese Arbeit adressiert zwei kritische Herausforderungen für zukünftige Fusionsreaktoren: (1) Schutz plasmazugewandter Komponenten durch kontrollierte Strahlung, und (2) Erreichen stabiler Ablösung bei hohen Strahlungsfraktionen (f_{\text{rad}}$\geq$85\%)
\item Drei Hauptbeiträge: Implementierung des ersten bolometer-basierten Echtzeit-Feedbacks am W7-X, Optimierung der Kanalauswahl durch systematische Sensitivitätsanalyse, und Validierung der MFR-Tomographie mit radial abhängiger Anisotropie-Gewichtung
\end{itemize}

\section{INHALTSVERZEICHNIS}

\begin{itemize}
\item Präsentationsstruktur: (1) Motivation - Fusionsenergie und W7-X-Herausforderungen, (2) Bolometer-Diagnostiksystem - 94 Kanäle mit $\sim$\SI{\1}{cm} räumlicher Auflösung, (3) Echtzeit-Feedback - erreichte f_{\text{rad}}$\geq$85\% mit Faktor $\geq$2 Wärmelast-Reduktion, (4) LOS-Sensitivität - identifizierte optimale Kanalkonfigurationen, (5) Tomographie - MFR mit RDA für 2D-Strahlungsrekonstruktion, (6) Schlussfolgerungen und zukünftige Arbeiten
\end{itemize}

\section{MOTIVATION}

\subsection{Kernfusion: Energieherausforderung}

\begin{itemize}
\item Kernfusion durch D-T-Reaktionen bietet saubere, reichlich vorhandene Energie - Deuterium aus Meerwasser, Tritium gezüchtet aus Lithium
\item Lawson-Kriterium: n_e·\tau_E·T $\geq$ 3$\times$10²¹ keV·s/m³ definiert Bedingungen für Netto-Energiegewinn (Q>1)
\item Erfordert extreme Bedingungen: T$\sim$10-\SI{\1}{keV} (100-200 Millionen °C) und Einschlusszeit \tau_E$\sim$1-\SI{\1}{s}
\item Magnetischer Einschluss nutzt starke Felder (\SI{\1}{T} am W7-X), um heißes Plasma von Materialwänden fernzuhalten
\item Stellaratoren wie W7-X erreichen Dauerstrichbetrieb durch 3D-geformte Spulen, im Gegensatz zu gepulsten Tokamaks
\end{itemize}

\subsection{Wendelstein 7-X}

\begin{itemize}
\item W7-X ist der weltweit größte und fortschrittlichste Stellarator: R=\SI{\1}{m}, a=\SI{\1}{m}, B=\SI{\1}{T}, optimiert für Dauerstrichbetrieb
\item Hauptherausforderung: Management der Wärmelasten auf plasmazugewandten Komponenten während langer Pulse (bis zu 30 Minuten geplant)
\item Inseldivertor-Konzept: 5 helikale Inselketten fangen Wärmefluss ab, Verbindungslängen 100-\SI{\1}{m} (10-100$\times$ länger als Tokamaks)
\item Ablösungsstrategie: Leistung in Abschälschicht und Rand abstrahlen, um Divertor-Wärmelasten von $\sim$\SI{\1}{MW}/m² auf <\SI{\1}{MW}/m² zu reduzieren
\item Ziel: stabile Ablösung bei f_{\text{rad}}$\geq$85-90\% ohne Plasmadisruption erreichen, während Kernleistung erhalten bleibt
\item Dies erfordert Echtzeitkontrolle der Strahlungsverteilung - Motivation für Bolometer-Feedbacksystem
\end{itemize}

\subsection{Plasmastrahlung}

\begin{itemize}
\item Plasmastrahlungsmechanismen: (1) Bremsstrahlung (frei-frei-Übergänge, $\sim$T_e^(1/2)), (2) Linienstrahlung (gebunden-gebunden, temperaturabhängig), (3) Rekombination (frei-gebunden)
\item Verunreinigungseinspeisung-Strategie: Injektion von niedrig-Z (N₂) oder hoch-Z (Ne, Ar) Gasen zur Verstärkung der Randstrahlung bei Vermeidung von Kernkontamination
\item Kohlenstoff und Sauerstoff sind intrinsische Verunreinigungen aus Wandwechselwirkungen - Kohlenstoff dominiert um Faktor 10²$\times$ über Sauerstoff am W7-X
\item Strahlungskühlungsfunktion P_{\text{rad}}/(V·n_e·n_imp) hat Maxima bei unterschiedlichen Temperaturen für jede Verunreinigung: C Maximum $\sim$10-\SI{\1}{eV} (Rand), O Maximum $\sim$\SI{\1}{eV}
\item Transport am W7-X: neoklassische Basislinie mit anomalen Beiträgen von Turbulenz (ITG, TEM-Moden)
\item Anomaler Transport erhöht Diffusion um Faktor 10-100 über neoklassisch, kritisch für Verunreinigungstransport-Modellierung
\item Strahlungsfraktion f_{\text{rad}} = P_{\text{rad}}/P_{\text{ECRH}} muss ausbalanciert sein: f_{\text{rad}}<70\% unzureichend für Divertorschutz, f_{\text{rad}}>95\% riskiert Kernkühlung und Einschlussdegradation
\item Optimales Betriebsfenster: f_{\text{rad}}=85-90\% erreicht Ablösung bei Erhalt der Plasmastabilität
\end{itemize}

\section{KERN-BOLOMETER AM W7-X}

\subsection{Bolometer-Diagnostik (Hardware)}

\begin{itemize}
\item Bolometer messen totale Plasmastrahlungsleistung (breitbandig, 0,2-\SI{\1}{nm}) durch resistive Heizung von Metallabsorbern
\item Metallwiderstand-Design: dünner Goldfilm (\SI{\1}{\micro m}) auf Glimmersubstrat, absorbiert Strahlung $\rightarrow$ Temperaturerhöhung $\rightarrow$ Widerstandsänderung
\item Wheatstone-Brücken-Konfiguration: AC-Anregung (\SI{\1}{V}, kHz) eliminiert 1/f-Rauschen, differentielle Messung verdoppelt Signal
\item Drei Kamera-Arrays bieten umfassende Plasmaabdeckung: HBC (32 Kanäle, horizontal), VBCl/r (2$\times$31 Kanäle, vertikal)
\item Insgesamt 94 Kanäle installiert, 61 funktional in OP1.2b-Kampagne (einige beschädigt durch ECRH-Streustrahlung)
\item Jeder Kanal integriert Strahlung entlang der Sichtlinie - nutzt Geometriematrix T und Etendue K_M zur Rekonstruktion der 2D-Verteilung
\item Räumliche Auflösung $\sim$\SI{\1}{cm} an magnetischer Achse, bestimmt durch Etendue-Breite und Beobachtungsgeometrie
\end{itemize}

\subsection{Leistung}

\begin{itemize}
\item Kalibrierung ist kritisch für absolute Leistungsmessungen - drei Methoden verwendet: (1) Laser, (2) elektrische ohmsche Heizung, (3) in-situ während Betrieb
\item In-situ ohmsche Heizungskalibrierung: U_{\text{cal}}=1,2-\SI{\1}{V} für \SI{\1}{s} anlegen, Stromevolution I(t) messen, R_M, \tau_M, \kappa_M aus exponentiellem Abfall extrahieren
\item Kalibrierungsparameter (OP1.2b Median): R_M=0,987$\pm$\SI{\1}{\ohm}, \kappa_M=0,724$\pm$0,055 A², \tau_M=110,57$\pm$\SI{\1}{m}s, bemerkenswert stabil über 1182 Experimente
\item Etendue K_M quantifiziert geometrische Lichtsammeleffizienz: berücksichtigt Detektor-Apertur-Geometrie, partielle Abschattung, winkelabhängige Transmission
\item Zeitauflösung konfigurierbar: {0,4, 0,8, 1,6, 3,2, 6,4, 12,8} ms Abtastung verfügbar, \SI{\1}{m}s für Feedback verwendet (balanciert Geschwindigkeit vs Rauschen)
\item Leistungsmetriken: SNR>1000 bei hoher Strahlung, Rauschen \sigma_ΔU=0,339$\pm$\SI{\1}{\micro V}, Gauß-Fehler \SI{\1}{\micro W}, Drift 0,057$\pm$\SI{\1}{\micro V}/s
\item Sehnen-Helligkeit P_{\text{ch}} = P_M/K_M liefert radiale Profilinformation - zeigt Strahlungsverteilung von Kern bis SOL
\item Absorptionseffizienz nahe Einheit (>95\%) über 0,2-\SI{\1}{nm} Wellenlängenbereich - wahrhaft breitbandige Messung
\end{itemize}

\section{ECHTZEIT-FEEDBACK}

\subsection{Inseldivertor-Konfiguration}

\begin{itemize}
\item W7-X nutzt Inseldivertor-Konzept: 5 helikale Inselketten (m/n=5/5 Moden) bilden natürliche Divertor-Ziele, 10 Divertor-Module insgesamt
\item Verbindungslängen L_c=100-\SI{\1}{m} (10-100$\times$ länger als Tokamaks) ermöglichen effiziente Strahlungskühlung entlang Feldlinien
\item Standardkonfiguration: Rotationstransformation \iota/2\pi=1 an magnetischer Achse, optimiert für Einschluss
\item Gaseinspeisung: 10 Piezoventile für präzise Kontrolle, Heliumstrahlventile für schnelle Injektion (Antwortzeit $\sim$\SI{\1}{m}s)
\item X-Punkte sind Strahlungs-Hotspots: Faktor 2-5$\times$ höhere Emissivität aufgrund von Verunreinigungsakkumulation und niedrigen Temperaturen
\item Bolometer-Kameras strategisch auf X-Punkte ausgerichtet für optimale Sensitivität gegenüber Ablösungsübergängen
\end{itemize}

\subsection{System-Implementierung}

\begin{itemize}
\item Erste bolometer-basierte Feedback-Implementierung am W7-X (OP1.2b-Kampagne, 2018) - vorherige Systeme nutzten Interferometrie
\item Hardware: NI 6321 DAQ-Karte (16-bit, 250 kS/s), PID-Regler in LabVIEW implementiert, parallele Verarbeitung für niedrige Latenz
\item Zwei Vorhersage-Proxies entwickelt: P_{\text{pred}}^(1) geometriegewichtet (nutzt mehrere Kanäle mit räumlichen Gewichten), P_{\text{pred}}^(2) Einzelkanal (einfachste Implementierung)
\item Optimale Kanalauswahl: VBC S={61,65,73,78,86} erreicht $\geq$85\% Vorhersagegenauigkeit für P_{\text{ECRH}}, beobachtet Separatrix/SOL-Region
\item Alternative Konfiguration: HBC S={4,7,8,20,23,27} für horizontale Abdeckung, nützlich wenn VBC-Kanäle ausfallen
\item PID-Parameter: K_p=0,5-2,0 (proportionale Verstärkung), K_i=0,1-0,5 (integrale Verstärkung), K_d=0 (keine derivative Aktion - zu rauschempfindlich)
\item Sollwert: f_{\text{rad}}^target=85-90\% balanciert Divertorschutz (erfordert $\geq$85\%) gegen Kernkühlung-Risiko (vermeidet >95\%)
\end{itemize}

\subsection{Latenz-Analyse}

\begin{itemize}
\item Akquisitionslatenz: \SI{\1}{m}s minimal (\SI{\1}{m}s Abtastung + \SI{\1}{m}s Verarbeitung), gemessen durch Laserdioden-Validierung
\item FIFO-Glättung: $\sim$\SI{\1}{m}s zusätzliche Verzögerung (M=10 Samples $\times$ \SI{\1}{m}s Abtastperiode) reduziert Rauschen
\item Gesamtschleifenlatenz: 150-\SI{\1}{m}s (plasmadominiert durch Gasventil-Antwort $\sim$50-\SI{\1}{m}s und Verunreinigungstransport $\sim$100-\SI{\1}{m}s)
\item Abtastrate-Kompromiss: \SI{\1}{m}s optimal balanciert Geschwindigkeit (schnelle Antwort) vs Rauschen (höhere Abtastraten verstärken Rauschen)
\item Laserdioden-Validierung: \SI{\1}{m}s Akquisitionslatenz experimentell bestätigt durch kontrollierte Lichtpulse
\item Zukünftige Verbesserungen: <\SI{\1}{m}s Akquisition möglich mit Hardware-Upgrade, prädiktive Algorithmen könnten Gesamtlatenz auf $\sim$50-\SI{\1}{m}s reduzieren
\end{itemize}

\subsection{XP20181010.32 Erfolge}

\begin{itemize}
\item Durchbruch-Experiment: Erste stabile Ablösung mit Wasserstoff-Einspeisung (vorherige Versuche nutzten Neon/Stickstoff)
\item Dauer: \SI{\1}{s} Feedback-Kontrolle, P_{\text{ECRH}}=\SI{\1}{MW}, Gesamtenergie \SI{\1}{MJ} in Plasma deponiert
\item Strahlungsfraktion: f_{\text{rad}}$\geq$85\% anhaltend aufrechterhalten, Spitzen bis 100\% während Ablösungsübergängen
\item Wärmelast-Reduktion: Faktor $\geq$2$\times$ erreicht (von $\sim$\SI{\1}{MW}/m² auf <\SI{\1}{MW}/m²), kritisch für Divertor-Lebensdauer
\item C²⁺-Spektroskopie-Signatur: Erscheint ab f_{\text{rad}}$\sim$50\%, zeigt Beginn der Ablösung, vollständige Ablösung bei f_{\text{rad}}>90\%
\item Vorhersagegenauigkeit: P_{\text{pred}}^(1) verfolgt P_{\text{rad}} innerhalb 10-15\% über gesamte Entladung, validiert Proxy-Konzept
\item Plasmaenergie: W_{\text{dia}}$\sim$0,8-\SI{\1}{MJ} trotz hoher Strahlung aufrechterhalten, zeigt Kernleistung erhalten bleibt
\item Reaktorrelevanz: Demonstriert Szenario anwendbar auf ITER/DEMO - hohe Strahlung ohne Einschlussdegradation
\end{itemize}

\subsection{Diskussion}

\begin{itemize}
\item \textbf{\1} Direkte P_{\text{rad}}-Messung (keine Modellannahmen), breitbandig (alle Wellenlängen), schnell (\SI{\1}{m}s Zeitauflösung)
\item \textbf{\1} Unabhängig von Verunreinigungsmix (misst totale Strahlung), funktioniert mit C, N₂, Ne, Ar-Einspeisung
\item \textbf{\1} Kanäle unterscheiden Kern vs Rand vs SOL-Strahlung, ermöglicht gezielte Kontrolle
\item \textbf{\1} Strahlung IST der Ablösungsmechanismus - direkteste mögliche Messung für Feedback
\item \textbf{\1} 50-\SI{\1}{m}s Ventil-Antwortzeit optimal für Plasma-Zeitskalen, schneller würde Rauschen verstärken
\item \textbf{\1} Parameter entladungsspezifisch, komplex zu optimieren, erfordert Erfahrung oder ML-Ansätze
\item \textbf{\1} 150-\SI{\1}{m}s Gesamtschleife begrenzt Bandbreite auf $\sim$2-\SI{\1}{Hz}, ausreichend für Ablösungskontrolle
\item \textbf{\1} <\SI{\1}{m}s Akquisition möglich, prädiktive Kontrolle könnte Latenz kompensieren, ML für automatisches PID-Tuning
\item \textbf{\1} Bolometer-Feedback ist direkteste verfügbare Methode für Strahlungskontrolle - keine Zwischenmodelle erforderlich
\item \textbf{\1} Konzept direkt anwendbar auf ITER, DEMO, und zukünftige Stellaratoren - universelle Lösung
\end{itemize}

\section{LOS-SENSITIVITÄT}

\subsection{LOS-Sensitivität (Motivation)}

\begin{itemize}
\item Zentrale Frage: Welche Bolometer-Kanäle sind am sensitivsten für Plasmaparameteränderungen und warum?
\item Ziel: Optimale Kanalauswahl für (1) Echtzeit-Feedback-Kontrolle und (2) Tomographie-Rekonstruktion identifizieren
\item Datensatz: 45 Entladungspaare aus OP1.2b, Parameterbereich P_{\text{ECRH}}=3,5-\SI{\1}{MW}, f_{\text{rad}}=0,3-1,0, verschiedene Verunreinigungen (N₂, Ne)
\item Sechs Evaluierungsmetriken entwickelt, um verschiedene Aspekte der Kanalsensitivität zu quantifizieren
\item Systematischer Ansatz: Kanal-für-Kanal-Vergleich über alle Entladungspaare, räumliche Kartierung der Sensitivität
\end{itemize}

\subsection{Methodik}

\begin{itemize}
\item Sechs Metriken definiert: (1) Korrelation \rho_c (lineare Beziehung), (2) gewichtete Abweichung \sigma_w (zeitgewichteter Fehler), (3) mittlere Abweichung \sigma_m (durchschnittlicher Fehler)
\item Fortsetzung: (4) FFT-Korrelation \rho_FFT (Frequenzbereichsähnlichkeit), (5) Fisher-Information I_F (Informationsgehalt), (6) \chi^2 (Anpassungsgüte)
\item Analyse-Workflow: Für jedes Entladungspaar, berechne alle Metriken für jeden Kanal, aggregiere über Paare, identifiziere konsistente Muster
\item Räumliche Kartierung: Sensitivität vs. \rho_pol (normalisierter poloidaler Fluss) zeigt, wo Kanäle beobachten
\item Optimale Kanalanzahl: Systematische Variation m=1-20 zeigt m=3-5 ausreichend für Feedback (Kompromiss zwischen Redundanz und Komplexität)
\item Hauptergebnis: Separatrix/SOL-beobachtende Kanäle (\rho_pol$\sim$0,95-1,05) zeigen höchste Sensitivität über alle Metriken
\end{itemize}

\subsection{Ergebnisse}

\begin{itemize}
\item Beste Kanäle identifiziert: VBC S={61,65,73,78,86} (beobachten Separatrix/SOL-Region), HBC S={4,7,8,20,23,27} (Rand-Fokus)
\item Sensitivität korreliert stark mit Strahlungsgradienten - Kanäle, die Regionen mit steilen $\partial$P_{\text{rad}}/$\partial$ρ beobachten, sind am sensitivsten
\item Kernkanäle (\rho_pol<0,8) zeigen niedrigere Sensitivität aufgrund flacher Strahlungsprofile - weniger Information über Plasmaänderungen
\item Randkanäle (\rho_pol$\sim$0,95-1,05) optimal für Feedback - erfassen Verunreinigungstransport und Strahlungsänderungen zuerst
\item Validierung: Experimentelle Feedback-Leistung mit ausgewählten Kanälen bestätigt Vorhersagen - stabile f_{\text{rad}}-Kontrolle erreicht
\end{itemize}

\section{TOMOGRAPHIE}

\subsection{Tomographie (Motivation)}

\begin{itemize}
\item Ziel: 2D-Strahlungsverteilung P_{\text{rad}}(R,Z) aus 1D-Sehnen-Helligkeitsmessungen rekonstruieren
\item Herausforderung: Unterbestimmtes inverses Problem - 61 Messungen, $\sim$4500 Gitterzellen, unendlich viele Lösungen möglich
\item Lösung: Minimum Fisher Regularization (MFR) mit radial abhängiger Anisotropie-Gewichtung (RDA)
\item MFR-Prinzip: Maximiere Entropie (Glätte) bei Erfüllung von Datenkonstraints - bevorzugt physikalisch plausible Lösungen
\item Validierung durch (1) synthetische Phantome mit bekannter Wahrheit und (2) experimentelle Daten mit unabhängigen P_{\text{rad}}-Messungen
\end{itemize}

\subsection{MFR mit RDA}

\begin{itemize}
\item Gitter-Konfiguration: 30$\times$150 Zellen (radial$\times$poloidal), Domäne 1,3²V_P umfasst Plasma + SOL + Divertor
\item Anisotropie-Gewichtung: k_{\text{ani}}(ρ) variiert radial - {2,0, 0,3} Standard (Kern anisotrop, Rand isotrop), {5,0, 0,5} hochanisotrop (starke Flussflächenausrichtung)
\item Physikalische Motivation: Strahlung folgt Flussflächen im Kern (anisotrop), wird isotrop im SOL/Divertor (kurze Verbindungslängen)
\item Konvergenz: N_T=10-20 Iterationen ausreichend für \chi^2<1,5, typisch $\sim$15 Iterationen für experimentelle Daten
\item Regularisierung: Glättung entlang Flussflächen im Kern, isotrope Glättung im Rand - balanciert Physik und Datenkonstraints
\item Implementierung: Python mit NumPy/SciPy, optimiert für Geschwindigkeit ($\sim$\SI{\1}{s} pro Rekonstruktion), parallelisierbar für Echtzeit
\end{itemize}

\subsection{Validierung}

\begin{itemize}
\item Phantom-Benchmarks: \SI{\1}{s}ynthetische Tests mit verschiedenen Strahlungsmustern (Kern-dominiert, Rand-dominiert, gemischt)
\item Phantom-Leistung: Rekonstruktionsqualität \rho_c>0,85 (Korrelation mit Wahrheit), \chi^2<2,5 (Anpassungsgüte), typisch \rho_c$\sim$0,90-0,95
\item Experimentelle Tests: 4 Entladungen aus OP1.2b, P_2D (integriert über Gitter) innerhalb 10-20\% von P_{\text{rad}} (unabhängige Messung)
\item Vergleich mit anderen Methoden: MFR+RDA übertrifft Standard-MFR (isotrop) und Tikhonov-Regularisierung - bessere räumliche Auflösung und Genauigkeit
\item Räumliche Auflösung: $\sim$\SI{\1}{cm} an magnetischer Achse, begrenzt durch Bolometer-Kanalgeometrie und Etendue-Breite
\item Robustheit: Stabil gegen Rauschen (SNR>100), fehlende Kanäle (bis zu 30\% Ausfall), und Kalibrierungsfehler ($\pm$10\%)
\item $\sim$2-\SI{\1}{s}/Frame CPU, GPU könnte <\SI{\1}{m}s
\item Zukunft: Echtzeit-Tomographie möglich
\end{itemize}

\section{SCHLUSSFOLGERUNGEN}

\subsection{Haupterfolge}

\begin{itemize}
\item Drei Hauptbeiträge bilden ein integriertes System für W7-X-Strahlungsmanagement:
\item (1) Echtzeit-Feedback: Erste bolometer-basierte Implementierung am W7-X, erreichte f_{\text{rad}}$\geq$85\% mit Faktor $\geq$2$\times$ Divertor-Wärmelast-Reduktion
\item (2) LOS-Optimierung: Systematische Analyse von 45 Entladungspaaren identifizierte optimale Kanalauswahl, m=3-5 Kanäle ausreichend für robuste Kontrolle
\item (3) Tomographie: MFR+RDA-Methode validiert durch 59 Phantom-Tests und 4 experimentelle Entladungen, liefert 2D-Strahlungsverteilung
\item Integration: Feedback-System nutzt optimierte Kanäle aus LOS-Analyse, Tomographie validiert räumliche Strahlungsverteilung
\item Reaktorrelevanz: Direkt anwendbar auf ITER und DEMO für Divertorschutz und Langpuls-Operationen
\end{itemize}

\subsection{Zusammenfassung}

\begin{itemize}
\item \textbf{\1} \SI{\1}{m}s Akquisitionslatenz (minimal), 150-\SI{\1}{m}s Gesamtschleife, PID-Regler mit dual-Proxy-Vorhersage
\item \textbf{\1} XP20181010.32 - \SI{\1}{s} Dauer, P_{\text{ECRH}}=\SI{\1}{MW}, f_{\text{rad}}$\geq$85\% aufrechterhalten, Faktor $\geq$2$\times$ Wärmelast-Reduktion, keine Disruption
\item \textbf{\1} 94 Kanäle installiert (61 funktional in OP1.2b), Kalibrierungsstabilität <5\% über 1182 Experimente, SNR>1000
\item \textbf{\1} Separatrix/SOL-beobachtende Kanäle (\rho_pol$\sim$0,95-1,05) optimal, VBC übertrifft HBC für Randsensitivität, mehrere Kanal-Sets erreichen $\geq$85\% Genauigkeit
\item \textbf{\1} C²⁺-Signatur erscheint bei f_{\text{rad}}$\sim$50\%, vollständige Ablösung bei f_{\text{rad}}>90\%, kritische Region bei \rho_pol$\approx$0,95
\item \textbf{\1} Kohlenstoff dominiert um Faktor 10²$\times$ über Sauerstoff, C³⁺/C²⁺-Verhältnis kritisch an LCFS für Strahlungsverteilung
\item \textbf{\1} Rekonstruktionen erreichen \rho_c>0,85, \chi^2<2,5, P_2D innerhalb 10-20\% von P_{\text{rad}}, zeigen X-Punkt-Lokalisierung und Inselketten-Struktur
\item \textbf{\1} P_2D^(Kern) stabil für quasi-stationäre Phasen, Standard-MFR unterschätzt um 15-25\% (SOL-Beitrag fehlt)
\item \textbf{\1} OP1.2b-Kampagne 2018, 1182 Experimente analysiert, 45 Entladungspaare für LOS-Studie, 4 für Tomographie-Validierung
\end{itemize}

\section{AUSBLICK}

\subsection{Zukünftige Arbeiten}

\begin{itemize}
\item \textbf{\1} (1) Akquisitionssystem-Upgrade auf <\SI{\1}{m}s Latenz für schnellere Feedback-Antwort, (2) Erweiterte Kanalabdeckung für bessere räumliche Auflösung
\item \textbf{\1} Prädiktive Algorithmen zur Antizipation von Strahlungsänderungen, maschinelles Lernen für adaptive Kontrolle
\item \textbf{\1} (1) PID-Parameter-Optimierung für verschiedene Plasmaregime, (2) automatische Kanalauswahl basierend auf Plasmazustand, (3) Gas-Injektions-Timing-Optimierung
\item \textbf{\1} STRAHL+EMC3-EIRENE 3D-Kopplung für vollständige Verunreinigungstransport-Simulation, Validierung gegen OP2.1-Daten
\item \textbf{\1} Verfeinerung von Zwei/Drei-Kammer-Modellen für Verunreinigungstransport, Entwicklung von Skalierungsgesetzen für verschiedene Betriebsregime
\item \textbf{\1} Erweiterte Phantom-Bibliothek (>\SI{\1}{T}ests), GPU-Beschleunigung für <\SI{\1}{m}s/Frame-Rekonstruktion
\item \textbf{\1} Integration in Feedback-System für räumliche Strahlungskontrolle, MARFE-Früherkennung und -Vermeidung
\item \textbf{\1} Bayes'sche Optimierung oder Gittersuche für k_{\text{ani}}(ρ)-Profil, adaptive Anisotropie basierend auf Plasmazustand
\item \textbf{\1} Direkt übertragbare Methoden, P_fus$\sim$\SI{\1}{MW}, f_{\text{rad}}$\geq$95\% erforderlich, q_div<\SI{\1}{MW}/m² Ziel
\item \textbf{\1} Anwendbar auf LHD (Stellarator), ITER, DEMO, SPARC (Tokamaks) - universelle Strahlungskontroll-Strategie
\item \textbf{\1} Vollautomatisches, adaptives, prädiktives Strahlungsmanagement-System für zukünftige Fusionsreaktoren
\end{itemize}

\section{ABSCHLUSS}

\subsection{Danksagung}

\begin{itemize}
\item Vielen Dank für Ihre Aufmerksamkeit - "Sin é, a chairde" (Irisch-Gälisch: "Das ist es, Freunde")
\item Fragen sind willkommen zu allen Aspekten: Feedback-Implementierung, LOS-Sensitivitätsanalyse, Tomographie-Methodik, experimentelle Ergebnisse
\item Diskussionsthemen: (1) Latenzreduktions-Strategien, (2) alternative Regularisierungsmethoden, (3) Parameter-Optimierungsansätze, (4) ITER/DEMO-Anwendbarkeit
\item Danksagung: W7-X-Betriebsteam, IPP-Greifswald-Kollegen, EUROfusion-Konsortium, Helmholtz-Gemeinschaft, Universität Greifswald
\end{itemize}

\section{BACKUP-FOLIEN - Übersicht}

\begin{itemize}
\item Backup-Folien liefern detaillierte technische Information für tiefgehende Fragen
\item Abgedeckte Themen: Plasmaphysik, STRAHL-Modellierung, Diagnostik-Details, LOS-Sensitivitätsmetriken, Tomographie-Vergleiche, DAQ-Latenz
\item Verfügbar, um spezifische Fragen zu Methodik, Implementierungsdetails oder alternativen Ansätzen zu adressieren
\end{itemize}

\vspace{1em}
\noindent\textit{Ende des Hauptpräsentations-Leitfadens}

\section{PLASMAPHYSIK}

\subsection{Plasmatransport am W7-X}

\begin{itemize}
\item Drei Transportmechanismen: klassisch, neoklassisch, anomal
\item W7-X-Optimierung für niedrige Kollisionalitätsregime
\item Anomaler Transport dominiert um Größenordnung
\item Verunreinigungstransport ladungsabhängig
\end{itemize}

\subsection{Verunreinigungstransport und Strahlung}

\begin{itemize}
\item Anomale Diffusion D_{\text{an}}$\sim$0,3-\SI{\1}{m}²/s dominiert
\item Neoklassische Konvektion erzeugt Einwärts-Pinch für hoch-Z
\item DEMO-Ziel: f_{\text{rad}}>\SI{\1}{m}it 70/30 Kern/SOL-Aufteilung
\item Kontrollherausforderung: Balance Schutz vs Einschluss
\end{itemize}

\subsection{Plasmaverunreinigungen}

\begin{itemize}
\item Kohlenstoff dominant: 10²$\times$ höher als Sauerstoff
\item Ionisationsstufen variieren mit Temperatur
\item Strahlungsverteilung entwickelt sich mit f_{\text{rad}}
\item Kern/SOL-Umkehr bei hoher f_{\text{rad}}
\end{itemize}

\subsection{Verunreinigungseffekte und Transportsensitivität}

\begin{itemize}
\item Brennstoffverdünnung beeinflusst Lawson-Kriterium
\item Diffusionssensitivität: 120\% Spitzenanstieg mit niedrigerem D
\item Separatrix-Profil kritisch für Ablösung
\item STRAHL validiert experimentelle Trends
\end{itemize}

\subsection{Plasmaablösung}

\begin{itemize}
\item Schwelle: T_e < \SI{\1}{eV} am Target
\item Volumetrische Verluste dominieren
\item Strahlungsverstärkungs-Rückkopplungsschleife
\item W7-X erreichte f_{\text{rad}}=100\% Ablösung
\end{itemize}

\section{MODELLIERUNG & ANALYSE}

\subsection{Verunreinigungseinspeisung-Modelle}

\begin{itemize}
\item Zwei-Kammer- und Drei-Kammer-Modelle
\item Erfassen grundlegende Zeitskalen: Wand, Plasma, SOL
\item R²>0,85 Korrelation mit Experimenten
\item Zukunft: Integration mit 3D-Codes
\end{itemize}

\subsection{Plasma-Leistungsbilanz}

\begin{itemize}
\item Energieerhaltung: P_bal $\approx$ 0
\item Kombinierte Unsicherheiten $\pm$1-\SI{\1}{MW}
\item Gefäßwärmelast 15-25\% von P_{\text{rad}}
\item Tomographie verbessert Stabilität
\end{itemize}

\section{DIAGNOSTIK-DETAILS}

\subsection{Detektorsensitivität: Etendue}

\begin{itemize}
\item K̃_M quantifiziert Lichtsammeleffizienz
\item Einheiten: mm⁻³ (inverse Volumen)
\item Räumliche Auflösung $\sim$4-\SI{\1}{cm} an Achse
\item Kritisch für absolute Leistungsmessungen
\end{itemize}

\subsection{Temperaturdrift-Einfluss}

\begin{itemize}
\item Hauptsystematischer Fehler in langen Pulsen
\item Graukörperstrahlung skaliert mit T⁴
\item Lineare Korrektur für langsame Drifts
\item Zukunft: temperaturabhängige Kalibrierung
\end{itemize}

\subsection{In-Situ-Kalibrierung}

\begin{itemize}
\item Essentiell wegen $\pm$5-10\% Variationen
\item Drei Parameter: R_M, \tau_M, \kappa_M
\item OP1.2b: außergewöhnliche Stabilität über 1182 Experimente
\item Kabelkorrektur verhindert systematische Fehler
\end{itemize}

\subsection{Detektorleistungs-Statistiken}

\begin{itemize}
\item SNR>1000 bei hoher Strahlung
\item Rauschen: 0,339$\pm$\SI{\1}{\micro V}
\item Minimale Degradation beobachtet
\item Erfüllt alle fusionsgerechten Anforderungen
\end{itemize}

\subsection{Bolometer-Gleichungs-Ableitung}

\begin{itemize}
\item AC-Anregung entfernt 1/f-Rauschen
\item Wheatstone-Brücke: ΔU ∝ ΔR_M
\item Leistungsbilanz: Heizung vs Kühlung
\item Trade-off: τ und κ invers verwandt
\end{itemize}

\subsection{Messalgorithmus}

\begin{itemize}
\item Drei Stufen: Initialisierung, Kalibrierung, Messung
\item Echtzeit-Feedback via NI 6321 DAQ
\item FIFO-Puffer verhindert Datenverlust
\item Latenzoptimierung: \SI{\1}{m}s Minimum
\end{itemize}

\section{STRAHL-MODELLIERUNG}

\subsection{STRAHL-Grundlagen}

\begin{itemize}
\item 1D-Verunreinigungstransport-Code
\item Flussflächen-gemittelte Größen
\item Zeitabhängige Evolution
\item Löst Kontinuität für jedes Z
\end{itemize}

\subsection{STRAHL-Verunreinigungstransport-Modellierung}

\begin{itemize}
\item Leitende Gleichung mit D und v
\item Kohlenstoffdominanz validiert
\item Kritische Ladungszustände: C³⁺/C²⁺
\item Strahlungsverschiebung mit f_{\text{rad}}
\end{itemize}

\subsection{STRAHL-Parametervariationen}

\begin{itemize}
\item Systematische Sensitivitätsstudie
\item Kernstrahlung stabil $\pm$10\%
\item Rand variiert $\pm$30-50\%
\item C³⁺/C²⁺ am empfindlichsten
\end{itemize}

\subsection{STRAHL: Strahlungsfraktions-Abhängigkeit}

\begin{itemize}
\item f_{\text{rad}}=90\% vs 100\% Vergleich
\item Spitze verschiebt sich von SOL zu LCFS
\item Kern/SOL-Verhältnis ändert sich 0,3$\rightarrow$0,7
\item Validiert Feedback-Kanalauswahl
\end{itemize}

\section{LOS-SENSITIVITÄTS-DETAILS}

\subsection{Kamerageometrie-Sensitivität}

\begin{itemize}
\item Aperturpositionierungs-Unsicherheiten
\item Vorwärtsmodellierung <1\% Unterschied
\item Fertigungstoleranzen akzeptabel
\item Validiert Ingenieurdesign
\end{itemize}

\subsection{Segmentierung und Sichtlinien-Kegel}

\begin{itemize}
\item N=8 optimal für Genauigkeit vs Berechnung
\item Voller Kegel vs infinitesimale Projektion
\item Richtige Segmentierung essentiell
\item Jedes Segment gewichtet durch Etendue
\end{itemize}

\subsection{Tomographie: RDA vs RGS Vergleich}

\begin{itemize}
\item RDA: absolute Differenzen, erhält Gradienten
\item RGS: relative Gradienten, verstärkt Merkmale
\item RDA gewählt für Robustheit
\item RGS nützlich für Post-Shot-Analyse
\end{itemize}

\subsection{Tomographie: Künstliche Kameras}

\begin{itemize}
\item VBCm und MIRh konzeptuelle Designs
\item 4-Kamera verbessert um 15-20\%
\item Trade-off: Kosten vs Verbesserung
\item Aktuelle 3-Kamera ausreichend
\end{itemize}

\subsection{Phantom-Tomographie: Grenzfälle}

\begin{itemize}
\item Invertierte Anisotropie testet Grenzen
\item Homogene Verteilung validiert keine Artefakte
\item Grenzfälle offenbaren Beschränkungen
\item Leiten Betriebsparameter
\end{itemize}

\subsection{LOS-Sensitivitätsbewertung (Kreuzkorrelation)}

\begin{itemize}
\item Zeitliche Ähnlichkeitsmetrik
\item Große Fehlerbalken: $\pm$40\%
\item X-Punkt-Kanäle am schlechtesten
\item Kritisch für Echtzeitkontrolle
\end{itemize}

\subsection{LOS-Sensitivitätsbewertung (Mittlere Abweichung)}

\begin{itemize}
\item Varianz-basierte Qualität
\item Fehlerbalken 10$\times$ größer als gewichtet
\item Kernkanäle konsistent
\item Rand optimal für Feedback
\end{itemize}

\subsection{LOS-Sensitivitätsbewertung (FFT-Korrelation)}

\begin{itemize}
\item Spektralanalyse-Ansatz
\item Identifiziert oszillatorisches Verhalten
\item Ergänzt Zeitbereich
\item Leitet Reglerentwurf
\end{itemize}

\subsection{LOS-Sensitivität: Plasmaparameter}

\begin{itemize}
\item P_{\text{ECRH}}-Korrelation: höher besser
\item Optimale f_{\text{rad}}: 0,4-0,75
\item T_e$\sim$\SI{\1}{keV} ideal
\item Validiert hohe f_{\text{rad}}-Szenarien
\end{itemize}

\subsection{DAQ-Latenz}

\begin{itemize}
\item Laser-basierte Charakterisierung
\item Abtastzeiten: 0,8-\SI{\1}{m}s getestet
\item Minimum: \SI{\1}{m}s erreicht
\item Totale Schleife: 150-\SI{\1}{m}s ausreichend
\end{itemize}

\vspace{2em}
\noindent\textbf{Dokument für PhD-Verteidigung vorbereitet}\\
\noindent\textbf{Gesamt-Abschnitte: 60+ Frames mit umfassenden Stichpunkten}\\
\noindent\textbf{Nutzen Sie diesen Leitfaden für konsistente Botschaften und technische Genauigkeit während der gesamten Präsentation}

\end{document}
